\documentclass[10pt,a4paper]{amsart}
\usepackage[margin=19mm]{geometry}
\usepackage{amsmath,amssymb,mathtools}
\usepackage[scr=rsfs,cal=euler]{mathalfa}
\usepackage{xfrac}
\usepackage[unicode,hidelinks,bookmarks=false]{hyperref}
\newcommand{\ie}{i.e.\ }
\newcommand{\cf}{cf.\ }
\newcommand{\resp}{respectively\ }
\newcommand{\Subsection}{\subsection}
\providecommand{\nobreakdash}{\nobreak\hbox{-}}
\setlength{\emergencystretch}{2em}
\multlinegap0pt
\usepackage{booktabs}
\newcommand{\theHalgorithm}{\arabic{algorithm}}
\usepackage{bbm,paralist,mathtools,tikz,nicefrac,url}
\newcommand{\ubar}[1]{\underaccent{\bar}{#1}}
\newcommand{\expect}[1]{\mathbb{E}\left[ #1 \right]} 
\newcommand{\cexpect}[2]{\mathbb{E}\left[ #1 \middle\vert #2 \right]} 
\newcommand{\mexpect}[2]{\mathbb{E}_{#1}\left[ #2\right]} 
\newcommand{\mcexpect}[3]{\mathbb{E}_{#1}\left[ #2 \middle\vert #3 \right]} 
\newcommand{\prob}[1]{\mathbb{P}\left( #1 \right)} 
\newcommand{\cprob}[2]{\mathbb{P}\left( #1 \middle\vert #2 \right)} 
\newcommand{\mprob}[2]{\mathbb{P}_{#1}\left( #2\right)} 
\newcommand{\mcprob}[3]{\mathbb{P}_{#1}\left( #2 \middle\vert #3 \right)} 
\newcommand{\abs}[1]{\left\vert #1 \right\vert}
\newcommand{\norm}[1]{\left\Vert #1 \right\Vert}
\newcommand{\oo}[1]{\left( #1 \right)}
\newcommand{\oc}[1]{\left( #1 \right]}
\newcommand{\co}[1]{\left[ #1 \right)}
\newcommand{\cc}[1]{\left[ #1 \right]}
\newtheorem{theorem}{Theorem}
\newtheorem{lemma}{Lemma}
\newtheorem{proposition}{Proposition}
\newtheorem{definition}{Definition}
\newtheorem{corollary}{Corollary}
\begin{document}
\title[Full Conformal Prediction under Stochastic Non-Conformity Me]{Full Conformal Prediction under Stochastic Non-Conformity Measure}
\author{Thanawat Sornwanee}
\date{}
\maketitle
\noindent\emph{Source and licence.} Full Conformal Prediction under Stochastic Non-Conformity Measure. Version arXiv:2606.28730v1 (CC BY 4.0). This material is distributed under the Creative Commons Attribution 4.0 International licence, http://creativecommons.org/licenses/by/4.0/, as recorded in the arXiv OAI licence metadata for this exact version and shown on the abstract page https://arxiv.org/abs/2606.28730v1. Full rights notice: licenses/arxiv-2606.28730-CC-BY-4.0.txt.\par\smallskip
\noindent\textit{Editorial scope.} Counted unit: the original \texttt{lemma} environment \texttt{lemma:exchangepvalue} at source lines 370--384 together with its complete proof at lines 385--394; the delivered document keeps source lines 356--401 so that every referenced label and every citation resolves inside it. Native editable source is the paper's own concatenated TeX; the AMS wrapper is editorial formatting only. No publisher class, upstream script or unrelated artwork is redistributed.
\section{p-Value \& permutation test}
\label{appendix:pperm}

    We define a valid p-value, and provide background on permutation test, which is the backbone conformal prediction.
    
    \begin{definition}
        [Valid p-Value]
        A random variable $P$ is a valid p-value if
        \begin{align*}
            \prob{P \le \alpha} \le \alpha 
        \end{align*}
        for all $\alpha \ge 0$.
    \end{definition}
    

    \begin{lemma}
    \label{lemma:exchangepvalue}
        Let $S:=\cc{S_i}_{i=1}^n$ be an exchangeable random variable where $S_i \in \mathcal{S}$ for each $i \in \{1,2,\dots, n\}$ and $\oo{\mathcal{S}, \succcurlyeq}$ is a total preorder. Define
        \begin{align*}
            P_1 := \frac{\sum_{i=1}^n \mathbf{1}_{S_1 \succcurlyeq S_i}}{n}.
        \end{align*}
        Then its distribution satisfies
        \begin{align*}
            \mathcal{L}(P_1)
            \succcurlyeq_{\mathrm{FOSD}}
            \mathrm{Unif}\!\left(\left\{\frac{1}{n},\ldots,1\right\}\right)
            \succcurlyeq_{\mathrm{FOSD}}
            \mathrm{Unif}([0,1]).
        \end{align*}
    \end{lemma}

    \begin{proof}
        Define $P_i:= \frac{\sum_{j=1}^n \mathbf{1}_{S_i \succcurlyeq S_j}}{n}$. Exchangeability of $S$ implies exchangeability of $P$. Note that, for any integer $k \in \{1,2,\dots, n\}$, there can be at most $k$ values of $i \in \{1,2,\dots, n\}$ where $P_i \le \frac{k}{n}$. Hence,
        \begin{align*}
            \prob{P_1 \le k/n}
            &= \expect{\frac{1}{n}
            \sum_{i=1}^n \mathbf{1}_{P_i \le k/n}} \\
            &\le \frac{k}{n},
        \end{align*}
        implying the first order stochastic dominations.
    \end{proof}

    \begin{corollary}
    \label{corollary:validpexchange}
        $P_1$ in the lemma~\ref{lemma:exchangepvalue} is a valid p-value if $S$ is exchangeable.
    \end{corollary}
    \begin{proof}
        From the lemma~\ref{lemma:exchangepvalue}, we have that $\mathbb{P}\left(P_1 \le \alpha\right) \le \left(\text{Unif}([0,1])\right)([0, \alpha]) = \alpha$, so $P_1$ is a valid p-value.
    \end{proof}

\end{document}
